% Contact sheet for the 23 Oct best-flight review figures (revised 23 Oct on Wesley's answers: colour, ls3c, r302_3,
% the spin column): every figure on its own page, at print size (the thesis's text width), with its one-line message.
% Build: pdflatex contact_sheet.tex (in this folder, or with \graphicspath).
\documentclass[12pt,a4paper]{article}
\usepackage[margin=1in]{geometry}
\usepackage{graphicx}
\usepackage{amsmath}
\usepackage{xcolor}
\graphicspath{{./}{/home/wesley/multi_drone_control/docs/figures/review_2026-10-23_best/}}
\pagestyle{plain}
\setlength{\parindent}{0pt}
\newcommand{\fig}[4]{%
  \section*{#1}
  \textit{#2}\par\smallskip
  {\small\color{black!60}#3}\par\vspace{1.2em}
  \includegraphics{#4}%
  \clearpage}
\begin{document}

\fig{1. The figure-8 with 2, 3 and 4 carriers (rig)}
{With two, three or four carriers the ring follows the 2\,m by 1\,m figure-8 within about 2\,cm (25, 19 and 17\,mm
RMS, clock-free), so at print scale its path lies on the reference.}
{The best rig figure-8 of each fleet size, each in its fleet colour (2 orange, 3 rose, 4 purple, throughout the set).
Four carriers flew on 7 Oct (the carry planner) at half the speed of the 22 Oct eights: to be replaced by Wed 28.}
{fig_traj.pdf}

\fig{2. Tracking error across fleet size and motion (recommended)}
{On the rig the ring keeps within 14--36\,mm of its path on every motion flown with 2, 3 or 4 carriers, and within
22\,mm of its hover point; the spin is the hardest (45 and 86\,mm). The twin (1\,m cables) keeps within 1--10\,mm in
every cell, with no trend across 2--4 carriers.}
{Each digit is a fleet size at its best flight's error, in its colour, on a log axis. Rig cells are not one variable
(day, code, speed differ); the pair hover and the four-carrier cells are to be replaced by Wed 28.}
{fig_summary.pdf}

\fig{2b. The same as a grid (alternative)}
{The same numbers as cells, darker = larger on a log scale, one scale for rig and twin.}
{A table in figure form: the thesis can carry it as a table instead.}
{fig_summary_grid.pdf}

\fig{3. The error's parts (recommended)}
{On the rig the sideways part of the error is the larger on every path but the spinny line (13--25\,mm against
6--26\,mm of height), and the ring runs 0.05--0.23\,s behind its schedule, reported apart.}
{Every chosen path flight but the spin (off this scale), in its fleet colour; open markers are the twin's spinny lines
(no rig flight with 2 or 4 carriers).}
{fig_components.pdf}

\fig{3b. The error along one lap (alternative)}
{Three carriers on the figure-8: sideways (+ inside the curve), height and lag against the ring's position, the four
tight bends shaded.}
{Shows where the error builds; busy for a results figure.}
{fig_components_eight.pdf}

\fig{4. Ring tilt (recommended)}
{With two carriers the ring tilts $7^\circ$ on average (to $14^\circ$) about the pair's line, a rotation no cable
resists; with three or four it stays under $6^\circ$ (median about $2^\circ$).}
{Boxes: quartiles; whiskers 5--95\,\%; dot: the peak. The same three figure-8 flights.}
{fig_tilt_box.pdf}

\fig{4b. Ring tilt along the eight (alternative)}
{The pair rocks at about 0.7\,Hz between $1^\circ$ and $14^\circ$; three and four stay low.}
{}
{fig_tilt.pdf}

\fig{5. Throttle against the cap (recommended)}
{No carrier reached its cap: four carriers flew at about 0.62 and three at 0.66 (cap 0.8); the pair flew at
0.72--0.75 with its cap raised to 0.9, where the default 0.8 would have left 0.01--0.05.}
{One row per carrier: the command sent to the flight controller over the figure-8.}
{fig_throttle_box.pdf}

\fig{5b. Throttle in time (alternative)}
{The same in time, one panel per fleet size.}
{Flat lines: it says less than 5.}
{fig_throttle.pdf}

\fig{6. Cable tension (recommended)}
{The planner's tensions sit on the static share of the ring's weight ($mg/n\sin 60^\circ$: 4.9, 3.2 and 2.4\,N)
within 3\,\%; the pull the rig measures reads 13--23\,\% above it.}
{Blue: each cable's planned tension; red: each carrier's measured pull (mocap acceleration less the modelled
thrust, so it carries any thrust-map error), its 5--95\,\% as a light bar. Means over the figure-8.}
{fig_tension.pdf}

\fig{6b. Cable tension in time (alternative)}
{Three carriers: the plan holds 3.3\,N while the measured pull wanders 3.4--4.8\,N.}
{}
{fig_tension_eight.pdf}

\fig{7. Reaction to a release (recommended)}
{When one of four carriers lets go unannounced, the planner sees the cable stretch past its 6\,cm threshold 0.42\,s
later, resizes to three 0.09\,s after that, and the ring's tilt, peaking at $6.5^\circ$, is back in its band 1.8\,s
after the release.}
{Rig, 8 Oct, hover. Released carrier red, survivors blue, the ring black. The band is the peak tilt of the 5\,s
before the release plus $0.5^\circ$.}
{fig_release.pdf}

\fig{7b. Reaction times on every rig release (companion)}
{On every rig release the planner detected the loose cable in 0.33--0.42\,s and resized 0.09--0.11\,s later; the
ring was back in band 1.8--3.8\,s after the release (tilt 6.5--17.8$^\circ$). The announced detach (r302\_3) resized
at once but tilted the ring to $15.4^\circ$.}
{Diamond: detected; bar: resized; grey line: to back in band; the number: peak tilt.}
{fig_reaction.pdf}

\fig{8. Planner solve time on the rig}
{Every solve finished inside the planner's 100\,ms period (median 16, 40 and 54\,ms); 16\,\% and 29\,\% of the ticks
passed the 60\,ms budget with 3 and 4 carriers.}
{Not one variable: four carriers is the 7 Oct carry planner. The controlled scaling is the S bench.}
{fig_solve.pdf}

\fig{9. Ring heading on the spinny line (rig)}
{Through the full $360^\circ$ turn the ring's heading runs about 0.5\,s behind its reference: the error holds near
$-12^\circ$ while the rate is constant (RMS $10.5^\circ$, at most $16^\circ$), about $2^\circ$ RMS once that delay is
taken out.}
{Three carriers, 22 Oct, today's constant-rate law (ls3c). The older law's flight (ls3, $229^\circ$): RMS $6.6^\circ$,
delay 0.44\,s.}
{fig_yaw.pdf}

\end{document}
